\section{Data Preprocessing}
\label{sec:solution:preprocessing}

The raw inputs released with the Ho Chi Minh City bus dataset arrive as
heterogeneous comma-separated value files with inconsistent column types,
misleading column labels in a small number of places, and no partitioning
suited to interactive analytical query. Before any analytical query can be
issued against the data, a preprocessing pipeline must normalise the inputs
to a uniform columnar storage format and register the result as queryable
views inside the analytical engine. The pipeline produces the input layer
consumed by the data architecture of
Section~\ref{sec:solution:architecture}.

The pipeline executes the following six steps in order:

\begin{enumerate}
  \item Ingest the raw comma-separated value files from the dataset release into staging memory frames, preserving every column as a string in the first pass to avoid silent coercion errors.
  \item Type-cast each column to the appropriate scalar or temporal type: integer for identifiers and counts, decimal for fare amounts, timestamp with timezone for event times, and date for calendar dates. Out-of-range values are written to a quarantine table for later inspection rather than dropped silently.
  \item Partition the two largest streams (the AFC fare records and the payments stream) by calendar day, producing one columnar file per day under the preprocessed data directory. The 90-day window thus yields 90 fare-record files and 90 payment files, which the analytical engine can prune on date predicates at query time.
  \item Rename misleading column labels. In particular, the raw \texttt{cash} field on the AFC stream is renamed to \texttt{stored\_value\_deduction} for the subset of records where the underlying value is a stored-value balance deduction made through a smartcard rather than an actual cash payment, so that downstream payment-channel breakdowns count each channel correctly.
  \item Tag the calendar windows that correspond to known natural-experiment periods: the Tết Lunar New Year window of 28 January 2026 to 5 February 2026 and the April 30 parade week (carried as a separate flag for the full-GPS phase). The tags are written as additional boolean columns on the calendar dimension produced by the preprocessing step.
  \item Register every resulting columnar file as a queryable view in the analytical engine, using a view-creation script that is regenerated each time the preprocessed directory changes. The view layer is the input to layer 1 of the warehouse, the layer of raw views described in Section~\ref{sec:solution:architecture}.
\end{enumerate}

The columnar storage format reduces total volume by approximately one order
of magnitude relative to the source comma-separated value files: the
17.88-million-row AFC stream compresses from approximately 4.3 gigabytes of
text to about 380 megabytes of columnar storage, and the payments stream
compresses by a similar factor. The reduction is the primary enabler of the
single-process deployability target stated in requirement N2 of
Section~\ref{sec:solution:problem:nonfunctional}. The registered view layer
that results from this section is the source of truth for every downstream
layer of the warehouse described in the next section.

\subsection{Validation checks}
\label{sec:solution:preprocessing:validation}

At the end of each pipeline run, four integrity checks are issued against
the freshly preprocessed view layer. A run that fails any of these checks
halts before the downstream warehouse build is triggered, so that the
warehouse never refreshes against partially preprocessed input.

\begin{itemize}
  \item Row-count regression per source. The row count of each source table is compared against the row count of the previous successful run; a deviation greater than 10\% in either direction halts the pipeline and raises an inspection request.
  \item Null-rate ceiling per column. The null rate of every column is computed and compared against a published per-column threshold; columns whose null rate exceeds the threshold halt the pipeline so that silent schema drift cannot reach the warehouse.
  \item Referential integrity for foreign keys. Every stop identifier and every route identifier on the fact streams is required to resolve to a row in the corresponding dimension; orphan identifiers are written to a quarantine table and counted, with a non-zero count halting the pipeline.
  \item Uniqueness of dimension primary keys. The primary key column of every dimension table is required to be unique and non-null; duplicates are written to the quarantine table and halt the pipeline.
\end{itemize}

The pipeline currently recomputes every preprocessed file from scratch on
every run. The trade-off considered is that scratch recomputation is
simple, fully reproducible, and free of state-management bugs, but is slow
on the full 50-day GPS volume; an incremental implementation would be
faster but would require a change-detection layer over the raw input
directory and a per-source watermark column. The simpler scratch policy is
used in the present phase; incremental processing is recorded as a
future-work item in Chapter~\ref{chapter:future}.
